% !TEX encoding = UTF-8
% !TEX program = pdflatex

\documentclass[
    version=last,      % Use the latest version of KOMA-Script
    paper=a4,          % Paper size A4
    fontsize=11pt,     % Font size
    english,           % Language settings for English
    fromalign=right,   % Sender address right-aligned
    parskip=half       % Half a line of space between paragraphs, no indent
]{scrlttr2}

\usepackage[utf8]{inputenc} % Allows direct input of UTF-8 characters
\usepackage[T1]{fontenc}    % Correct hyphenation for characters with accents
\usepackage{babel}          % Language-specific typography (date, etc.)
\usepackage[scaled]{helvet} % Use Helvetica as the sans-serif font
\renewcommand\familydefault{\sfdefault} % Set sans-serif as the default font family
\usepackage{lmodern}        % Modern font for better text quality
\usepackage{graphicx}       % For including graphics (e.g., a scanned signature)
\usepackage[normalem]{ulem} % Provides underline and strikethrough
\usepackage{enumitem}       % For customizing lists
\usepackage{hyperref}       % For creating hyperlinks in the PDF
\hypersetup{colorlinks=true, linkcolor=black, urlcolor=black}

% --- Sender Data ---
% This data is displayed in the letterhead
\setkomavar{fromname}{Your Name}
\setkomavar{fromaddress}{Your Street 123 \\ 12345 Your City}

% --- Date ---
% \today automatically inserts the current date.
% Alternatively, set a fixed date: \setkomavar{date}{June 28, 2025}
\setkomavar{date}{\today}

% --- Subject Line ---
\setkomavar{subject}{Subject of the Letter}

\begin{document}

% --- Recipient Address ---
% This will appear in the address field (suitable for windowed envelopes)
\begin{letter}{
    Recipient Name \\
    Recipient Street 456 \\
    54321 Recipient City
}

% --- Salutation ---
\opening{Dear Recipient,}

This is the main content of your letter. You can write multiple paragraphs here.

This is a new paragraph, separated by a blank line in the source code. The `parskip=half` option ensures there is some vertical space between paragraphs.

% --- Closing Formula ---
\closing{Sincerely,}

% --- Signature ---
% This leaves space for a handwritten signature.
% The name below it will be printed.
\setkomavar{signature}{[Your Name]}

\end{letter}
\end{document}
